\section{The Instrument: Method and Experimental Setup}
\label{sec:instrument}

\subsection{Task, inputs and metrics}
\label{sec:task}
Each model ranks a graph's layout configurations by runtime, using TPUGraphs' official train/validation
split for each collection. Models train with the ListMLE ranking loss~\cite{xia2008listwise} on up to
128 sampled configurations per graph per epoch; validation scores every configuration. Every model
receives the same static inputs: 140 features per node plus 18 layout features per configurable node.
These already include the convolution sizes (\texttt{shape\_dimensions} and \texttt{window\_size}) and
the kernel layout (layout features 12--17); we took the feature names from Google's featurizer source
and checked them against the data. The models therefore lack only the compiler's tiling and an
explicit MXU-utilisation number. We report Kendall's $\tau$, computed per graph and averaged across
graphs~\cite{kendall1938measure}, and Top-$K$ slowdown, the autotuner-facing metric: the best true
runtime among the model's top $K$ configurations divided by the best runtime overall.

\subsection{The three models}
\label{sec:models}
Each model adds one layer of structure while the parameter count stays roughly constant. The flat
baseline is GraphSAGE~\cite{hamilton2017inductive}: three SAGEConv layers with residual connections,
read out as a mean over all nodes plus a separate mean over configurable nodes, with hidden width 172
(236,157 parameters). The hierarchical model replaces flat message passing with two DiffPool pooling
stages~\cite{ying2018hierarchical}, the first clustering to at most 128 clusters (235,553 parameters).
The contention-edge model adds edge-aware GINE convolutions over the data-flow
edges~\cite{hu2019strategies} plus two edge types that link nearby operations competing for the matrix
unit (window of 5) and for memory bandwidth (window of 15) (234,120 parameters). These models are the
instruments being compared, not contributions.

\subsection{The protocol}
\label{sec:protocol}
Every comparison uses parameter-matched arms (within 0.3\% for flat and hierarchical, 0.9\% for
contention-edge), the same five seeds for every arm, and paired two-sided $t$-tests on
per-graph-averaged $\tau$ over held-out graphs, with Wilcoxon signed-rank tests as a direction check
and Benjamini--Hochberg correction across the three-arm family~\cite{benjamini1995controlling}. Each
control has a job. Seed-to-seed spread in $\tau$, about 0.04 to 0.07, is as large as the effects being
measured, so arms are paired by seed to cancel shared noise. Capacity differences fake wins: our
original ``DiffPool beats flat'' result came from 73\% more parameters (235,553 vs 136,321) and
vanished once the flat baseline was matched (Section~\ref{sec:xla-random}).

\subsection{Resolution of the design}
\label{sec:resolution}
With five seeds, the smallest detectable change in $\tau$ is 0.066--0.116 (Table~\ref{tab:mde}).
Effects below that are invisible to this design, not necessarily absent, and every null in this paper
is read against its minimum detectable effect (MDE). Detecting a change of 0.05 at a worst-case
$\sigma$ of 0.070 would need 18 seeds; because resolution improves only with the square root of the
seed count, the seed budgets common in this literature cannot resolve differences of the size papers
claim. The number of evaluation graphs matters less for the average: on xla:default, the bootstrap
standard deviation of $\tau$ is already stable at 7 graphs (median 0.0145, against 0.0131 at 20;
Appendix~C). It matters more for the 3-graph convolution subset in Section~\ref{sec:capstone}, which
we test with leave-one-graph-out and held-out-family checks.

\begin{table}[t]
\caption{Minimum detectable change in Kendall's $\tau$ (paired $t$-test, $\alpha = 0.05$, power 0.8)
by number of seeds $n$ and seed-to-seed standard deviation $\sigma$ of $\Delta\tau$.}
\label{tab:mde}
\centering
\small
\setlength{\tabcolsep}{4pt}
% Generated by scripts/paper_figures.py via power_analysis.min_detectable_effect
% (paired t, alpha=0.05, power=0.8)
\begin{tabular}{lcccccc}
\toprule
$\sigma$ (paired) & $n=3$ & $n=5$ & $n=7$ & $n=10$ & $n=15$ & $n=20$ \\
\midrule
0.040 & 0.124 & 0.066 & 0.051 & 0.040 & 0.031 & 0.026 \\
0.055 & 0.170 & 0.091 & 0.070 & 0.055 & 0.043 & 0.036 \\
0.070 & 0.217 & 0.116 & 0.089 & 0.070 & 0.054 & 0.046 \\
\bottomrule
\end{tabular}

\end{table}

\subsection{Evaluation choices}
\label{sec:eval-choices}
Two choices shape the numbers. First, all models use BatchNorm with batch statistics, so a
configuration's score depends on the evaluation batch it is scored in. $\tau$ comparisons use
evaluation batch 8; the Top-1 analyses in Sections~\ref{sec:result2} and~\ref{sec:capstone} use batch
32, their registered batch. Appendix~C gives every headline contrast at both, and the text notes each
case where the batch changes a verdict. Second, all arms train at one shared learning rate of 1e-3,
which matches them but does not tune them; Sections~\ref{sec:result1} and~\ref{sec:capstone} report the
pre-registered checks that tune each arm. The exception is nlp:default, whose amended runs use 1e-4 for
both arms (Section~\ref{sec:nlp-default}).
